% BASIC keyword reference.

\guidechapter{BASIC KEYWORD REFERENCE}{%
  \item Every wedge statement, function and constant
  \item Exact syntax, tokens, results and side effects
  \item One short example for every entry
}

\begin{referencechapter}

\section{Common Rules}

All entries require the BASIC wedge from Chapter~2. Custom statements work in
direct and program modes, but BASIC V2 cannot dispatch one immediately after
\bas{then}; branch to a line containing the statement instead. None of the
custom keywords has an abbreviated entry.

Statements report their SDK result through \bas{uerr}; target-specific status
is retained in \bas{udev}. Numeric syntax errors remain ordinary BASIC errors.
Dedicated filename statements retain at most 63 filename bytes. The raw
\bas{uci} statement retains 256 reply bytes, 40 status bytes and 64 argument
bytes.

\section{Raw Command Interface}

\basicapi{UCI}{statement}

\entryfield{Syntax}{\bas{uci} or \bas{uci} \textit{target, command}
[\bas{,} \textit{argument} \ldots]}
\entryfield{Modes}{Direct and program}
\entryfield{Token}{\cmd{\$CC} (204)}
\entryfield{Result}{\bas{uerr}, \bas{udev}, \bas{ulen}, \bas{ust\$},
\bas{udat\$} and \bas{ubyte(}}

\entryhead{Purpose}

Executes any published UCI command. A bare \bas{uci} aborts an unfinished
exchange and returns the interface to idle. Target and command are bytes.
String arguments contribute their bytes; numeric arguments normally contribute
one byte. Known commands automatically widen arguments required by their wire
format. \bas{uw(} and \bas{ul(} force widths explicitly.

\entryhead{Examples}

\begin{screen}
10 UCI UCTRL,1\\
20 PRINT UERR;ULEN;UDAT\$
\end{screen}

\begin{outputbox}
 0  14 CONTROL TARGET
\end{outputbox}

\begin{screen}
10 UCI UDOS1,38\\
20 UCI UDOS1,38,0\\
30 UCI UDOS1,240,UW(4660),UL(65536)
\end{screen}

\begin{resultbox}
The first two calls demonstrate an omitted and explicit optional argument. The
third sends 4660 as two bytes and 65536 as four, both least-significant byte
first.
\end{resultbox}

\basicapi{UERR}{numeric function}

\entryfield{Syntax}{\bas{uerr}}
\entryfield{Modes}{Direct and program}
\entryfield{Token}{\cmd{\$CD} (205)}
\entryfield{Returns}{The last \cmd{ULTIMATE\_*} result; zero is
\cmd{ULTIMATE\_OK}}

\entryhead{Purpose}

Tests whether the most recent wedge statement succeeded. It is initialized to
zero when the wedge starts. Appendix~A defines every result code.

\begin{screen}
10 ULOAD "FONT.PRG",49152\\
20 IF UERR THEN PRINT "LOAD FAILED";UERR
\end{screen}

\begin{resultbox}
The program continues silently when the load succeeds and prints the SDK result
when it fails.
\end{resultbox}

\basicapi{UDEV}{numeric function}

\entryfield{Syntax}{\bas{udev}}
\entryfield{Modes}{Direct and program}
\entryfield{Token}{\cmd{\$CE} (206)}
\entryfield{Returns}{The target's last raw status number, from 0 through 65535}

\entryhead{Purpose}

Adds firmware-specific detail to diagnostics. Do not use it for normal program
control; use \bas{uerr}, whose meanings are stable across targets.

\begin{screen}
10 UCI UDOS1,1\\
20 IF UERR THEN PRINT "DEVICE STATUS";UDEV
\end{screen}

\begin{resultbox}
UDEV is consulted only after UERR reports failure.
\end{resultbox}

\basicapi{ULEN}{numeric function}

\entryfield{Syntax}{\bas{ulen}}
\entryfield{Modes}{Direct and program}
\entryfield{Token}{\cmd{\$CF} (207)}
\entryfield{Returns}{The retained reply length, from 0 through 256}

\entryhead{Purpose}

Reports how many data bytes the most recent \bas{uci} command retained. It is
zero before the first command. A longer reply retains the first 256 bytes and
sets \bas{uerr} to \cmd{ULTIMATE\_ERR\_TRUNCATED}.

\begin{screen}
10 UCI UCTRL,1\\
20 FOR I=0 TO ULEN-1:PRINT UBYTE(I);:NEXT
\end{screen}

\begin{resultbox}
The loop visits exactly the retained reply bytes.
\end{resultbox}

\basicapi{UST\$}{string function}

\entryfield{Syntax}{\bas{ust\$}}
\entryfield{Modes}{Direct and program}
\entryfield{Token}{\cmd{\$D0} (208)}
\entryfield{Returns}{The status bytes from the most recent \bas{uci} command as
a BASIC string}

\entryhead{Purpose}

Exposes the target's status text exactly as received, up to the wedge's
40-byte status buffer. Use \bas{uerr} and \bas{udev} when a stable numeric test
is needed.

\begin{screen}
10 UCI UCTRL,1\\
20 PRINT UST\$
\end{screen}

\begin{outputbox}
00,OK
\end{outputbox}

\basicapi{UDAT\$}{string function}

\entryfield{Syntax}{\bas{udat\$}}
\entryfield{Modes}{Direct and program}
\entryfield{Token}{\cmd{\$D1} (209)}
\entryfield{Returns}{Up to 255 retained reply bytes as a BASIC string}

\entryhead{Purpose}

Returns textual or binary reply data in one value. BASIC strings cannot contain
more than 255 bytes, so use \bas{ubyte(255)} to inspect the final byte of a
256-byte retained reply.

\begin{screen}
10 UCI UCTRL,1\\
20 PRINT UDAT\$
\end{screen}

\begin{outputbox}
CONTROL TARGET
\end{outputbox}

\basicapi{UBYTE(}{numeric function}

\entryfield{Syntax}{\bas{ubyte(}\textit{index}\bas{)}}
\entryfield{Modes}{Direct and program}
\entryfield{Token}{\cmd{\$D2} (210), including the opening parenthesis}
\entryfield{Returns}{Reply byte \textit{index}, or zero when the index is
outside the retained reply}

\entryhead{Purpose}

Reads binary replies without converting them to strings. Valid retained
indices are 0 through 255.

\begin{screen}
10 UCI UCTRL,1\\
20 PRINT UBYTE(0);UBYTE(1);UBYTE(999)
\end{screen}

\begin{outputbox}
 67  79  0
\end{outputbox}

\basicapi{UW(}{UCI argument modifier}

\entryfield{Syntax}{\bas{uw(}\textit{numeric expression}\bas{)}}
\entryfield{Valid only}{As an argument inside a \bas{uci} statement}
\entryfield{Token}{\cmd{\$D3} (211), including the opening parenthesis}
\entryfield{Produces}{Two little-endian argument bytes}

\entryhead{Purpose}

Overrides the raw command's normal argument width. It is useful for firmware
commands newer than the wedge's shape table.

\begin{screen}
UCI UDOS1,240,UW(4660)
\end{screen}

\begin{resultbox}
The argument bytes are 52,18 (\$34,\$12).
\end{resultbox}

\basicapi{UL(}{UCI argument modifier}

\entryfield{Syntax}{\bas{ul(}\textit{numeric expression}\bas{)}}
\entryfield{Valid only}{As an argument inside a \bas{uci} statement}
\entryfield{Token}{\cmd{\$D4} (212), including the opening parenthesis}
\entryfield{Produces}{Four little-endian argument bytes}

\entryhead{Purpose}

Forces a signed 32-bit BASIC value into the raw argument stream.

\begin{screen}
UCI UDOS1,240,UL(65536)
\end{screen}

\begin{resultbox}
The argument bytes are 0,0,1,0.
\end{resultbox}

\section{Target Constants}

\basicapi{UDOS1}{numeric constant}
\entryfield{Syntax}{\bas{udos1}}
\entryfield{Token}{\cmd{\$D5} (213)}
\entryfield{Returns}{1, the first Ultimate DOS target}
\begin{screen}
UCI UDOS1,1
\end{screen}
\begin{resultbox}
The IDENTIFY command is sent to the primary filesystem service.
\end{resultbox}

\basicapi{UDOS2}{numeric constant}
\entryfield{Syntax}{\bas{udos2}}
\entryfield{Token}{\cmd{\$D6} (214)}
\entryfield{Returns}{2, the second Ultimate DOS target}
\begin{screen}
UCI UDOS2,1
\end{screen}
\begin{resultbox}
The IDENTIFY command is sent to the optional second filesystem service.
\end{resultbox}

\basicapi{UNET}{numeric constant}
\entryfield{Syntax}{\bas{unet}}
\entryfield{Token}{\cmd{\$D7} (215)}
\entryfield{Returns}{3, the network target}
\begin{screen}
UCI UNET,1
\end{screen}
\begin{resultbox}
The IDENTIFY command is sent to the network service.
\end{resultbox}

\basicapi{UCTRL}{numeric constant}
\entryfield{Syntax}{\bas{uctrl}}
\entryfield{Token}{\cmd{\$D8} (216)}
\entryfield{Returns}{4, the machine and cartridge-control target}
\begin{screen}
UCI UCTRL,1
\end{screen}
\begin{resultbox}
The IDENTIFY command is sent to the control service.
\end{resultbox}

\basicapi{UIEC}{numeric constant}
\entryfield{Syntax}{\bas{uiec}}
\entryfield{Token}{\cmd{\$D9} (217)}
\entryfield{Returns}{5, the SoftwareIEC target}
\begin{screen}
UCI UIEC,1
\end{screen}
\begin{resultbox}
The IDENTIFY command is sent to the high-speed IEC service.
\end{resultbox}

\basicapi{UHTTP}{numeric constant}
\entryfield{Syntax}{\bas{uhttp}}
\entryfield{Token}{\cmd{\$DA} (218)}
\entryfield{Returns}{6, the HTTP target}
\begin{screen}
UCI UHTTP,1
\end{screen}
\begin{resultbox}
The IDENTIFY command is sent to the HTTP client service.
\end{resultbox}

\section{Turbo Control}

\basicapi{UTURBO}{statement and numeric function}

\entryfield{Syntax}{\bas{uturbo} \textit{index} or \bas{uturbo}}
\entryfield{Modes}{Direct and program}
\entryfield{Token}{\cmd{\$DB} (219)}
\entryfield{Statement result}{\bas{uerr}}
\entryfield{Function result}{Current speed index, or 255 when turbo registers
are unavailable}

\entryhead{Purpose}

Sets or reads the Ultimate~64 CPU-speed index. Indices 0, 1, 2 and 3 portably
mean 1, 2, 3 and 4\,MHz. Indices 4 through 15 are model-dependent. The
statement leaves the badline setting unchanged.

\begin{screen}
10 S=UTURBO:IF S=255 THEN PRINT "NO TURBO":END\\
20 UTURBO 0:UTURBO 1:UTURBO 2:UTURBO 3\\
30 UTURBO 15:IF UERR THEN END\\
40 PRINT UTURBO\\
50 UTURBO S
\end{screen}

\begin{outputbox}
 15
\end{outputbox}

The first four settings demonstrate the portable 1, 2, 3 and 4\,MHz speeds;
15 selects the model's maximum. Indices 4--14 are accepted but deliberately
not assigned a frequency here because their meanings differ by model.

\section{Files And Directories}

\basicapi{ULOAD}{statement}

\entryfield{Syntax}{\bas{uload} \textit{filename\$}
[\bas{,} \textit{address}]}
\entryfield{Modes}{Direct and program}
\entryfield{Token}{\cmd{\$DC} (220)}
\entryfield{Result}{\bas{uerr}; raw target status in \bas{udev}}

\entryhead{Purpose}

Loads a PRG from the current Ultimate directory. Its two-byte header is always
consumed. Omitting the address, or supplying zero, uses the header's address;
any other 16-bit value overrides it.

\begin{screen}
10 ULOAD "FONT.PRG"\\
20 IF UERR THEN PRINT "LOAD FAILED";UERR;UDEV\\
30 ULOAD "FONT.PRG",49152
\end{screen}

\begin{resultbox}
The first call uses the PRG's address. The second loads its body beginning at
\$C000. The caller must protect live memory at either destination.
\end{resultbox}

\basicapi{UBLOAD}{statement}

\entryfield{Syntax}{\bas{ubload} \textit{filename\$, address, length}}
\entryfield{Modes}{Direct and program}
\entryfield{Token}{\cmd{\$DD} (221)}
\entryfield{Result}{\bas{uerr}; raw target status in \bas{udev}}

\entryhead{Purpose}

Loads at most \textit{length} raw file bytes into a nonzero 16-bit C64 address.
It does not consume or interpret a PRG header. Address and length must both be
nonzero.

\begin{screen}
10 UBLOAD "FONT.BIN",49152,2048\\
20 IF UERR THEN PRINT "LOAD FAILED";UERR;UDEV
\end{screen}

\begin{resultbox}
At most 2048 bytes are written beginning at \$C000.
\end{resultbox}

\basicapi{USAVE}{statement}

\entryfield{Syntax}{\bas{usave} \textit{filename\$, start, length}}
\entryfield{Modes}{Direct and program}
\entryfield{Token}{\cmd{\$DE} (222)}
\entryfield{Result}{\bas{uerr}; raw target status in \bas{udev}}

\entryhead{Purpose}

Creates or replaces a file with \textit{length} bytes beginning at the 16-bit
C64 address \textit{start}. Length must be nonzero.

\begin{screen}
10 USAVE "SCREEN.BIN",1024,1000\\
20 IF UERR THEN PRINT "SAVE FAILED";UERR;UDEV
\end{screen}

\begin{resultbox}
SCREEN.BIN receives the 1000-byte text screen.
\end{resultbox}

\basicapi{UDIR}{statement}

\entryfield{Syntax}{\bas{udir}}
\entryfield{Modes}{Direct and program}
\entryfield{Token}{\cmd{\$DF} (223)}
\entryfield{Result}{\bas{uerr}; directory names are printed as they arrive}

\entryhead{Purpose}

Lists the current Ultimate directory, one entry per line. A slash follows each
directory. Normal end-of-directory leaves \bas{uerr} at zero.

\begin{screen}
UDIR
\end{screen}

\begin{outputbox}
TEMP/
\end{outputbox}

The output above is the controlled root directory used by the hardware test;
real directory contents vary.

\section{RAM Expansion}

\basicapi{USTASH}{statement}

\entryfield{Syntax}{\bas{ustash} \textit{C64 address, REU address, length}}
\entryfield{Modes}{Direct and program}
\entryfield{Token}{\cmd{\$E0} (224)}
\entryfield{Result}{\bas{uerr}}

\entryhead{Purpose}

Copies C64 memory into the RAM expansion with DMA. The C64 address is 16-bit,
the REU address is 24-bit, and the length is 16-bit; a zero length means 65536
bytes. The transfer is complete when the statement returns.

\begin{screen}
10 USTASH 1024,0,1000\\
20 IF UERR THEN PRINT "NO REU"
\end{screen}

\begin{resultbox}
The 1000-byte text screen is copied to REU address zero.
\end{resultbox}

\basicapi{UFETCH}{statement}

\entryfield{Syntax}{\bas{ufetch} \textit{C64 address, REU address, length}}
\entryfield{Modes}{Direct and program}
\entryfield{Token}{\cmd{\$E1} (225)}
\entryfield{Result}{\bas{uerr}}

\entryhead{Purpose}

Copies data from the RAM expansion into C64 memory. Address widths and the
zero-means-65536 length convention are the same as \bas{ustash}.

\begin{screen}
10 UFETCH 1024,0,1000\\
20 IF UERR THEN PRINT "NO REU"
\end{screen}

\begin{resultbox}
The first 1000 REU bytes replace the C64 text screen.
\end{resultbox}

\basicapi{UREU}{numeric function}

\entryfield{Syntax}{\bas{ureu}}
\entryfield{Modes}{Direct and program}
\entryfield{Token}{\cmd{\$E2} (226)}
\entryfield{Returns}{Installed REU size in 64K banks, or zero when unavailable}

\entryhead{Purpose}

Tests for a RAM expansion and measures it in one operation. Multiply by 64 for
kilobytes.

\begin{screen}
10 PRINT UREU\\
20 PRINT UREU*64
\end{screen}

\begin{outputbox}
 32\\
 2048
\end{outputbox}

This is the actual result with the hardware-test configuration set to a 2\,MB
REU; the value follows the machine's setting.

\evidence{\cmd{tools/gen\_keywords.py}; \cmd{src/basic/dispatch.s};
\cmd{tests/emulator/basic.suite}; \cmd{tests/hardware/basictest.py}}

\end{referencechapter}
